\section{Methods}
\label{sec:methods}

\subsection{Datasets and targets}
\label{sec:methods-datasets-targets}

The in-domain data are the aligned neutral tmQM transition-metal complexes with \(\omega\)B97M-V/def2-SVPD energies and fixed train, validation, and test molecule identifiers \cite{balcells2020tmqm,mardirossian2016wb97mv,rappoport2010def2svpd}. The full panel contains 56,830 training, 7,104 validation, and 7,104 test molecules. The six descriptor residuals use 56,821 training and 7,103 validation molecules, while all 13 primary targets share the same 7,104 test molecules. Both architectures use identical target values and split memberships within each target. DFTB success filtering leaves 56,540 training, 7,068 validation, and 7,068 test molecules; failed calculations were not imputed. QM9 is used only for OOD evaluation; its energies were recomputed at the same level with Psi4 to avoid conflating chemical-domain shift with electronic-structure mismatch \cite{ramakrishnan2014quantum,parrish2017psi4}.

For method \(m\), molecular energies were first corrected using fixed element-wise atomic references,
\begin{equation}
  E_{\rm ref}^{m}(i)=E_{\rm raw}^{m}(i)-\sum_Z n_Z(i)\epsilon_Z^{m},
  \label{eq:atomic-reference}
\end{equation}
then corrected by a linear model of elemental counts fitted only on the tmQM training split to obtain \(E_{\rm ref,BoE}^{m}\). Atomic references are fixed external inputs, not parameters fitted to held-out molecular energies. Separate composition models are fitted for DFT, xTB, and DFTB and frozen for validation, test, and QM9 evaluation. The DFT quantity is the parent target. Delta targets subtract either a fitted descriptor baseline or the consistently referenced xTB/DFTB energy. Full target-generation, masking, and molecule-ID alignment details are in SI Methods.

\subsection{Baseline models}
\label{sec:methods-baselines}

BoE denotes elemental-count features and the standard BoE reference is linear. A composition-only control fits a nonlinear MLP to the same counts and atom-referenced DFT energy, thereby changing baseline flexibility without adding structural information. BoB denotes the study-specific RDKit-derived typed-bond-count vector, not the conventional Coulomb-matrix Bag-of-Bonds descriptor. PRE expands local element-resolved pair distances in Gaussian radial basis functions up to a spatial cutoff, and PRE+A augments PRE with central-element-resolved triplet angle distributions. BoB, PRE, and PRE+A each have linear ridge and MLP fits to the DFT-BoE parent target; all four MLP baselines share hidden widths of 512 and 128 and the same training protocol. Ridge regularization was selected by five-fold training-only cross-validation, with preprocessing refitted within each fold. Final baseline parameters and preprocessing were fitted using only the corresponding tmQM training split. The validation split was used to select each MLP checkpoint, while test labels were used only for final evaluation. Baseline fitting seed 101 defines the canonical residual targets; fitting-seed sensitivity is reported separately in SI Note~9. GFN2-xTB and DFTB are fixed semiempirical baselines \cite{bannwarth2019gfn2xtb,elstner1998sccdftb,hourahine2020dftbplus}. DFTB uses SCC-DFTB/DFTB2 with the PTBP Slater--Koster set, without third-order or dispersion corrections \cite{dftborg2025ptbp}, and only paired successful calculations; failures were not imputed. Feature dimensions, hyperparameters, and failure audits are in SI Methods and SI Note~8.

\subsection{GNN training}
\label{sec:methods-gnn-training}

GemNet-T is the primary downstream model and SchNet is the architecture control \cite{gasteiger2021gemnet,schutt2023schnetpack2}. Models were trained on energies only, without forces or periodic boundary conditions, using fixed splits, a \(6.0~\text{\AA}\) cutoff, and a common FairChem/OCP-style protocol. Every target was standardized by its own training-set mean and population standard deviation. Within each architecture, architecture settings, loss, AdamW configuration, initial learning rate, scheduler rule, batch size, and checkpoint criterion were held fixed across targets. AdamW with AMSGrad used an initial learning rate of \(5\times10^{-4}\), MSE energy loss, and an 80-epoch budget; batch sizes were 16 for SchNet and four for GemNet. ReduceLROnPlateau responded independently to each validation trajectory, with factor 0.8 and patience three, producing target-specific learning-rate schedules. Each run's best-validation checkpoint was evaluated on both the tmQM test set and QM9. All 13 primary targets use independent GNN seeds 101, 102, and 103 for both architectures; the two DFTB targets use seed 101 only for both architectures, giving 82 runs in total. Reported primary results are means of seed-level metrics, not ensemble predictions, with sample standard deviations shown where applicable. Single-seed DFTB results have no across-seed uncertainty estimate. Exact seed-level results are reported in SI Note~9. Complete configurations and training diagnostics are in SI Methods and Supplementary Table~S3.

\subsection{Roughness and transfer diagnostics}
\label{sec:methods-roughness-diagnostics}

GemNet roughness uses standardized PRE+A features and SchNet roughness uses standardized PRE features. A cosine \(k=10\) nearest-neighbour graph was built on the 7,104-molecule tmQM test cohort. Zero-variance feature channels were removed, and retained channels were divided by the graph cohort's population standard deviation without centring. This unsupervised graph preprocessing is distinct from train-only baseline preprocessing. Unique undirected union edges were assigned similarity-based weights using the locally scaled Gaussian kernel in Equation~\ref{eq:dirichlet-edge-weight}.
\begin{equation}
  w_{ij}=\exp\!\left(-\frac{d_{ij}^{2}}{\sigma_i\sigma_j}\right),
  \label{eq:dirichlet-edge-weight}
\end{equation}
Here, \(d_{ij}\) is the cosine distance between standardized descriptor vectors for molecules \(i\) and \(j\), and \(\sigma_i\) is a local distance scale, defined as the distance from molecule \(i\) to its \(k\)-th nearest neighbour.

Using these edge weights, we quantify target variation through the IQR-normalized graph Dirichlet roughness, \(D_{\rm IQR}\) \cite{belkin2006manifold}:
\begin{equation}
  D_{\rm IQR}(y)=\frac{1}{\sum_{\{i,j\}\in\mathcal E}w_{ij}}
  \sum_{\{i,j\}\in\mathcal E}
  w_{ij}\frac{(y_i-y_j)^2}{\operatorname{IQR}_{\rm test}(y)^2}.
  \label{eq:dirichlet-roughness}
\end{equation}
Here, \(\mathcal E\) is the set of unique undirected graph edges. For targets with nonzero IQR, this formulation normalizes local squared target differences by the squared central target spread, making \(D_{\rm IQR}\) invariant to nonzero uniform target rescaling.

Both \(D_{\rm IQR}\) and \(n\text{-MAE}_{\rm test}\) use the corresponding test-target IQR; graphs and weights remain fixed across primary targets within each architecture. DFTB controls use separately constructed valid-subset test graphs and are excluded from primary regressions and pooled statistics.

Semi-log regressions use unweighted ordinary least squares across the 13 primary target means, with an intercept and \(\ln D_{\rm IQR}\) as predictor. Raw-error, raw-roughness, variance-normalization, partial-correlation, and leave-one-target-out sensitivity analyses are reported in SI Note~7. The same note reports 5,000 joint test-label permutations, which preserve each target's distribution and IQR while breaking molecular-neighbourhood alignment, and comparisons of target pairs whose IQRs differ by at most 10\%. Leave-one-target-out fits are influence checks, not cross-validated prediction scores. These test-set analyses are retrospective diagnostics performed after model selection; test labels were not used to fit baselines, schedule training, or select checkpoints.

\subsection{OOD transfer evaluation}
\label{sec:methods-transfer}

tmQM-trained models were evaluated without retraining on QM9 geometries, using frozen tmQM training-only baseline parameters. Residual predictions were combined with their corresponding baselines and mapped to the common DFT-BoE parent before full-target correlation and ranking metrics were calculated. Primary parent-energy comparisons use 128,988 common QM9 molecules and three GNN seeds. The DFTB residual is shown on its 128,978-molecule intersection with this cohort, alongside a matched-subset direct seed-101 reference; this matches evaluation molecules, not training cohorts. Mean-shifted MAE removes one scalar mean-error offset estimated separately for each predictor from the evaluated QM9 reference labels and is used only as a post hoc oracle diagnostic, not an independently calibrated prediction score. Figure~\ref{fig:qm9-transfer} reports GemNet raw and mean-shifted MAEs in Ha, without IQR normalization, and Kendall correlations; exact coverage, offset magnitudes, and additional metrics are reported in SI Note~5. Detailed reconstruction and evaluation procedures are in SI Note~5.
